\documentclass[12pt]{article}
\usepackage[utf8]{inputenc}
\usepackage{amsmath}
\usepackage{amssymb}
\usepackage{geometry}
\usepackage{listings}
\usepackage{xcolor}
\usepackage{hyperref}

\geometry{a4paper, margin=1in}

\definecolor{codegreen}{rgb}{0,0.6,0}
\definecolor{codegray}{rgb}{0.5,0.5,0.5}
\definecolor{codepurple}{rgb}{0.58,0,0.82}
\definecolor{backcolour}{rgb}{0.95,0.95,0.92}

\lstset{
    backgroundcolor=\color{backcolour},   
    commentstyle=\color{codegreen},
    keywordstyle=\color{magenta},
    numberstyle=	iny\color{codegray},
    stringstyle=\color{codepurple},
    basicstyle=	tfamily\small,
    breakatwhitespace=false,         
    breaklines=true,                 
    captionpos=b,                    
    keepspaces=true,                 
    numbers=left,                    
    numbersep=5pt,                  
    showspaces=false,                
    showstringspaces=false,
    showtabs=false,                  
    tabsize=2
}

	itle{Documentation: A* Search Algorithm for Grid Pathfinding}
\author{AI Lab - Assignment 3}
\date{	oday}

\begin{document}

\maketitle

\section{Introduction}
The A* search algorithm is a complete and optimal heuristic search algorithm used extensively in pathfinding and graph traversal. In this implementation (	exttt{grid\_solver\_v1.py}), the algorithm is applied to a 2D grid with obstacles and 8-directional movement.

\section{Mathematical Formulation}
A* evaluates nodes by combining the cost to reach the node and the estimated cost to reach the goal. The evaluation function is defined as:
\begin{equation}
    f(n) = g(n) + h(n)
\end{equation}
Where:
\begin{itemize}
    \item $g(n)$: The actual cost of the path from the start node to node $n$.
    \item $h(n)$: The heuristic cost estimate from node $n$ to the goal.
    \item $f(n)$: The estimated total cost of the path through node $n$.
\end{itemize}

\subsection{Heuristic Function}
The implementation uses the 	extbf{Euclidean Distance} as the heuristic function, which is admissible for 8-directional grids if the diagonal movement cost is scaled correctly:
\begin{equation}
    h(n) = \sqrt{(x_n - x_{goal})^2 + (y_n - y_{goal})^2}
\end{equation}

\subsection{Movement Costs}
In a grid with diagonal movement, the step cost $c(n, n')$ between adjacent nodes is:
\begin{equation}
    c(n, n') = 
    \begin{cases} 
    1.0 & 	ext{if move is horizontal or vertical} 
    \sqrt{2} \approx 1.414 & 	ext{if move is diagonal}
    \end{cases}
\end{equation}

\section{Algorithm Workflow}
The algorithm maintains two sets:
\begin{enumerate}
    \item 	extbf{Open Set}: A priority queue containing nodes to be evaluated, sorted by $f(n)$.
    \item 	extbf{G-Score Map}: Stores the lowest cost to reach each node found so far.
\end{enumerate}

At each step, the node with the minimum $f(n)$ is popped from the Open Set. For each valid neighbor, a tentative $g(n)$ is calculated. If this cost is lower than the previously recorded $g(n)$, the neighbor's scores are updated, and it is added to the Open Set.

\section{Python Implementation Mapping}
\begin{itemize}
    \item 	exttt{heuristic(a, b)}: Implements the Euclidean distance calculation.
    \item 	exttt{get\_neighbors(node)}: Logic for valid grid bounds and obstacle detection (	exttt{grid[r][c] == 0}).
    \item 	exttt{solve()}: The main loop utilizing 	exttt{heapq} for the priority queue.
\end{itemize}

\end{document}
